\paragraph{window-\texorpdfstring{$6$}{6} 的内禀自校准与边界证书：末位均值比率、四块谱数域、返时母函数与半群阈值}
本段在不引入外部标定的前提下，把若干与 window-\(6\) 锚点直接相连、且可由有限数据审计的结论压缩为一条可独立引用的链条：其一是 bin-fold 两态一致渐近所诱导的 \(\varphi\) 自校准统计量；其二是 boundary uplift 的四步平移同构所给出的家族数窗口量子化与最小窗唯一性；其三是四块粗粒化核的不可约三次谱数域与 \(S_3\) 伽罗瓦群；其四是边界块 \(U_R\) 的首返时间分布闭式与通量倒数律；其五是由 tail-offset 三元组 \(\{21,34,55\}\) 所生成数值半群的 Frobenius 阈值与由此导出的“维数筛法有限性”。
\par

\begin{theorem}[黄金比的内禀自校准：由末位条件均值比率恢复 \texorpdfstring{$\varphi$}{phi}]\label{thm:xi-foldbin-intrinsic-self-calibration-phi-lastbit-means}
在 bin-fold 两态一致渐近（定理 \ref{thm:fold-bin-two-state-asymptotic}）成立的口径下，
令 \(d_m^{\mathrm{bin}}(x):=\bigl|(\mathrm{Fold}^{\mathrm{bin}}_m)^{-1}(x)\bigr|\)（\(x\in X_m\)），并按末位划分
\[
X_m^{(0)}:=\{x\in X_m:\ x_m=0\},
\qquad
X_m^{(1)}:=\{x\in X_m:\ x_m=1\}.
\]
定义可观测统计量
\[
\widehat\varphi_m
\ :=\
\frac{\frac{1}{|X_m^{(0)}|}\sum_{x\in X_m^{(0)}} d_m^{\mathrm{bin}}(x)}{\frac{1}{|X_m^{(1)}|}\sum_{x\in X_m^{(1)}} d_m^{\mathrm{bin}}(x)}.
\]
则存在与 \(m\) 无关常数 \(C>0\) 使得
\[
\bigl|\widehat\varphi_m-\varphi\bigr|\ \le\ C\left(\frac{\varphi}{2}\right)^m,
\]
从而 \(\widehat\varphi_m\) 以指数精度收敛到 \(\varphi\)。
该结论给出一条无需外部标定的自校准机制：仅凭折叠纤维多重度与末位可见位即可恢复增长常数。
\end{theorem}

\begin{proof}
由定理 \ref{thm:fold-bin-two-state-asymptotic}，存在常数 \(C_0>0\) 使对充分大 \(m\) 与任意 \(x\in X_m\) 有一致估计
\[
\Bigl|d_m^{\mathrm{bin}}(x)-\varphi^{-x_m}\Bigl(\frac{2}{\varphi}\Bigr)^m\Bigr|\le C_0.
\]
对 \(X_m^{(0)}\)、\(X_m^{(1)}\) 分别取均值，得
\[
\frac{1}{|X_m^{(0)}|}\sum_{x\in X_m^{(0)}} d_m^{\mathrm{bin}}(x)
=\Bigl(\frac{2}{\varphi}\Bigr)^m+O(1),
\qquad
\frac{1}{|X_m^{(1)}|}\sum_{x\in X_m^{(1)}} d_m^{\mathrm{bin}}(x)
=\varphi^{-1}\Bigl(\frac{2}{\varphi}\Bigr)^m+O(1),
\]
其中 \(O(1)\) 的常数与 \(m\) 无关。
两式相除并用 \((2/\varphi)^{-m}=(\varphi/2)^m\) 得
\[
\widehat\varphi_m=\varphi\cdot\frac{1+O((\varphi/2)^m)}{1+O((\varphi/2)^m)}
=\varphi+O\!\left(\left(\frac{\varphi}{2}\right)^m\right),
\]
从而存在常数 \(C\) 满足所示误差界。证毕。
\end{proof}

\begin{theorem}[家族数的窗口量子化与最小窗唯一性：\texorpdfstring{$N_f=|X_m^{\mathrm{bdry}}|$}{Nf=|Xbdry|}]\label{thm:xi-family-number-window-quantization-minimal-window}
固定每代 \(SU(5)\) 手征计数 \(15\) 为外部输入，并要求 “\(SO(10)\)-型最小补全” 在同一窗口 \(m\) 上由边界 singlet 通道完成，即满足离散恒等式
\[
15N_f+|X_m^{\mathrm{bdry}}|=16N_f.
\]
则必有
\[
N_f=|X_m^{\mathrm{bdry}}|=F_{m-2},
\]
从而家族数在该制度下只能取 Fibonacci 值。
特别地，若要求 \(N_f=3\)，则唯一解为 \(m=6\)（最小窗唯一性），并与 window-\(6\) uplift 三分支中 \(F_9=34\) 分支的离散对齐一致（表 \ref{tab:window6_family_uplift_lock} 与 \eqref{eq:window6_family_singlet_completion}）。
\end{theorem}

\begin{proof}
由所给离散恒等式立得 \(N_f=|X_m^{\mathrm{bdry}}|\)。
另一方面，由定理 \ref{thm:boundary-shift4-uplift-isomorphism} 的 uplift 同构 \(X_{m-4}\simeq X_m^{\mathrm{bdry}}\)（\(m\ge 6\)）与 \(|X_n|=F_{n+2}\) 得
\(|X_m^{\mathrm{bdry}}|=|X_{m-4}|=F_{m-2}\)，从而得到 \(N_f=F_{m-2}\)。
若 \(N_f=3\)，则 \(F_{m-2}=3\) 强制 \(m-2=4\)，即 \(m=6\)，且该解唯一。证毕。
\end{proof}

\begin{theorem}[window-\(6\) 四块粗粒化核的谱数域刚性：不可约三次域与 \texorpdfstring{$S_3$}{S3} 伽罗瓦群]\label{thm:xi-window6-four-block-kernel-cubic-field-s3}
令四块骨架的粗粒化 Markov 核 \(T\) 如 \eqref{eq:window6_edge_flux_skeleton}，
并记其特征多项式分解为
\[
\chi_T(t)=(t-1)\cdot p(t),
\qquad
p(t)=\frac{48114t^3+7263t^2-506t-55}{48114}.
\]
则：
\begin{enumerate}
  \item 三次因子 \(p(t)\) 在 \(\mathbb Q[t]\) 上不可约；
  \item 其判别式 \(\Delta=\mathrm{Disc}(48114t^3+7263t^2-506t-55)\) 非平方（见 \eqref{eq:window6_edge_flux_skeleton_arithmetic}）；
  \item 因而分裂域的伽罗瓦群为 \(S_3\)，并且 \(T\) 的全部非平凡本征值均为代数次数 \(3\) 的实数。
\end{enumerate}
因此非平凡谱生成一个不可交换的三次数域 \(K_T=\mathbb Q(\lambda)\)，并且“粗粒化混合”所需的最小代数复杂度在数域意义下为三次且不可降阶。
\end{theorem}

\begin{proof}
对三次多项式，\(\mathbb Q[t]\) 上不可约当且仅当无有理根。
把整系数多项式
\(
f(t):=48114t^3+7263t^2-506t-55
\)
模 \(13\) 化简，得
\(
f(t)\equiv t^3+9t^2+t+10\pmod{13}
\)。
直接检验 \(t\in\mathbb F_{13}\) 的全部取值可知该三次多项式在 \(\mathbb F_{13}\) 上无根，故在 \(\mathbb F_{13}[t]\) 上不可约，从而 \(f\) 在 \(\mathbb Q[t]\) 上不可约，\(p\) 亦不可约。
\par
由 \eqref{eq:window6_edge_flux_skeleton_arithmetic}，判别式 \(\Delta\) 非平方。
对不可约三次多项式，其伽罗瓦群为 \(A_3\) 当且仅当判别式为平方；否则为 \(S_3\)。
因此本例伽罗瓦群为 \(S_3\)。
又因 \(\Delta>0\)，三根互异且均为实数，从而 \(T\) 的三个非平凡本征值均为代数次数 \(3\) 的实数。证毕。
\end{proof}

\begin{proposition}[边界块 \texorpdfstring{$U_R$}{UR} 的首返时间分布完全可解：有理型母函数与显式二阶矩]\label{prop:xi-window6-UR-first-return-time-closed-form}
以上述四状态核 \(T\) 为动力学（\eqref{eq:window6_edge_flux_skeleton}），令 \(\tau_R\) 为从状态 \(U_R\) 出发的首返时间（首次回到 \(U_R\) 的步数，取值于 \(\mathbb N\)）。
则其概率母函数 \(F_R(z):=\mathbb E[z^{\tau_R}]\) 为有理函数，并满足闭式证书
\input{sections/generated/eq_window6_ur_first_return_time_pgf}
从而严格矩量为
\[
\mathbb E[\tau_R]=\frac{32}{3},
\qquad
\mathrm{Var}(\tau_R)=\frac{1695268}{15417}.
\]
此外，\(F_R\) 的收敛半径由分母多项式的唯一正根 \(\rho_R>1\) 决定；因此 \(\tau_R\) 具有严格指数型尾部衰减。
\end{proposition}

\begin{corollary}[\texorpdfstring{$(1+3+8)$}{(1+3+8)} 的倒数通量：边界稳态通量常数恰为 \texorpdfstring{$1/\dim\mathfrak g_{\mathrm{SM}}$}{1/dim gSM}]\label{cor:xi-window6-reciprocal-flux-boson12}
在 window-\(6\) 的四块骨架证书中（\eqref{eq:window6_edge_flux_skeleton}），边界块 \(U_R\) 的割大小与稳态通量常数满足
\[
\Phi_R=\frac{|\partial U_R|}{64\cdot 6}=\frac{1}{12}.
\]
若将标准模型规范玻色子计数解释为 \(\dim\mathfrak g_{\mathrm{SM}}=1+3+8=12\)，则得到严格倒数律
\[
\Phi_R=\frac{1}{\dim\mathfrak g_{\mathrm{SM}}}.
\]
\end{corollary}

\begin{corollary}[复位事件时钟的分支频率刚性：\texorpdfstring{$\{34,55\}$}{\{34,55\}} 的极限频率由 \texorpdfstring{$\varphi^{-1}$}{phi^{-1}} 唯一决定]\label{cor:xi-window6-reset-clock-gap-frequencies}
对 window-\(6\) 的复位事件集合 \(\mathcal R_6\)（定义 \ref{def:terminal-reset-events}），其返回时间二值谱为
\(\{F_9,F_{10}\}=\{34,55\}\)（定理 \ref{thm:terminal-reset-events-sturmian}；并见表 \ref{tab:fold_zeck_reset_event_gaps_m6}）。
记字母 \(A:=55\)、\(B:=34\)，则差分字母序列为 Fibonacci 替换 \(\sigma(A)=AB,\sigma(B)=A\) 的不动点；因此两间隔的极限频率存在且为
\[
\lim_{N\to\infty}\frac{\#\{n\le N:\ r_{n+1}-r_n=55\}}{N}=\frac{1}{\varphi},
\qquad
\lim_{N\to\infty}\frac{\#\{n\le N:\ r_{n+1}-r_n=34\}}{N}=\frac{1}{\varphi^2}.
\]
从而 “\(SO(10)\) 尾部偏移 \(34\)” 在无限时钟尺度上获得不可调的概率权重，为统一分支的概率化选择给出可审计的确定性频率定律。
\end{corollary}

\begin{theorem}[尾部偏移半群的 Frobenius 阈值：\texorpdfstring{$\langle 21,34,55\rangle$}{<21,34,55>} 的导通界]\label{thm:xi-window6-tail-semigroup-frobenius-threshold}
在 window-\(6\) uplift 的离散尾部偏移口径下，非零偏移仅可能取自
\(\{21,34,55\}=\{F_8,F_9,F_{10}\}\)（见 \S\ref{subsubsec:window6-terminal-propositions} 的相应结论链）。
令数值半群
\[
S_6:=\langle 21,34,55\rangle\subset\mathbb N.
\]
则 \(S_6\) 为余有限半群，且其 Frobenius 数与导通数满足
\input{sections/generated/eq_window6_tail_semigroup_frobenius}
换言之，对任意整数 \(n\ge 660\) 都存在 \(a,b,c\in\mathbb N\) 使
\(
n=21a+34b+55c
\)，而 \(659\) 为最大不可表示数。
\end{theorem}

\begin{corollary}[维数筛法的有限性：在尾部可加假设下 \texorpdfstring{$D>671$}{D>671} 不能再由维数否证]\label{cor:xi-window6-dimension-sieve-finite}
若在“额外规范模态仅通过尾部偏移的可加叠加进入”这一可审计假设下，将可达玻色子总数集合写为
\[
\mathcal D_{\mathrm{reach}}:=12+S_6,
\]
则由定理 \ref{thm:xi-window6-tail-semigroup-frobenius-threshold} 的导通数立得：
所有 \(D\ge 12+660=672\) 均属于 \(\mathcal D_{\mathrm{reach}}\)，而不可达维数集合 \(\mathbb N\setminus\mathcal D_{\mathrm{reach}}\) 为有限集，且其最大元为 \(12+659=671\)。
因此，超过 \(671\) 的维数不能再仅由尾部可达性被否证；必须引入更细的签名 \(\Sigma\)、退化谱、或谱穿越证书等更高分辨率不变量。
\end{corollary}

\begin{proposition}[从半群阈值到可证伪预测：硬否证与软筛选的分层]\label{prop:xi-window6-falsifiability-layering-from-semigroup}
在区间 \(D\le 671\) 内，\(\mathcal D_{\mathrm{reach}}\) 之外的维数点构成绝对否证点：任何声称由 window-\(6\) 尾部可加机制产生的统一候选若具有这些维数，则与协议的可达性约束直接矛盾。
相反，对 \(\mathcal D_{\mathrm{reach}}\cap[0,671]\) 内的维数点，维数本身不再构成否证，必须进一步满足 Zeckendorf--Fibonacci 分辨率签名的一致审计（例如禁相邻约束、边界层激活的无冗余性以及谱/PSD 证书的兼容性）。
因此可证伪性在该框架中被精确分割为：低维有限区间内由半群阈值提供的硬否证，与同区间内由签名/谱证书提供的软筛选两类。
\end{proposition}

